% ============================================================
% CAPITOLUL 8 — EVALUAREA SOFTWARE, PERFORMANȚEI ȘI SCALABILITĂȚII
% ============================================================

\chapter{Evaluarea Software, Performanței și Scalabilității}
\label{ch:evaluare}

\section{Scopul evaluării}
\label{sec:scop_evaluare_software}

Evaluarea software completează analiza numerică din
Capitolul~\ref{ch:comparatie_calibrare}. Dacă analiza numerică măsoară cât de
bine se potrivește un model cu datele, evaluarea software măsoară cât costă
obținerea acelui rezultat într-un sistem real. În aplicația implementată,
simularea nu este un apel local către o funcție, ci un flux distribuit: API-ul
creează o rulare, Redis transportă job-ul, runner-ul execută solver-ul,
YugabyteDB persistă rezultatele, iar frontend-ul afișează statusul și graficul.

Prin urmare, evaluarea trebuie să separe patru componente:
\begin{itemize}
    \item costul numeric al solverului;
    \item costul de serializare și persistare a rezultatelor;
    \item costul operațional introdus de coadă și de baza de date;
    \item costul de utilizare perceput în interfața web.
\end{itemize}

Această separare este importantă deoarece optimizarea unei singure componente nu
garantează îmbunătățirea întregului sistem. Un solver rapid poate fi mascat de
scrieri lente în baza de date, iar un model cu puține compartimente poate deveni
costisitor dacă este rulat de multe ori în cadrul calibrării.

\section{Metrici colectate}
\label{subsec:metrici}

\begin{table*}[!htbp]
\centering
\caption{Metricile colectate sau derivabile pentru evaluarea software.}
\label{tab:metrici_benchmark}
\begin{tabularx}{0.95\textwidth}{l l X}
\toprule
\textbf{Metrică} & \textbf{Nivel} & \textbf{Interpretare} \\
\midrule
\texttt{wall\_time\_ms} & runner & durata observată între începutul execuției și finalizarea persistării rezultatelor \\
\texttt{cpu\_time\_ms} & runner & timpul de procesor consumat de procesul numeric \\
\texttt{peak\_memory\_mb} & runner & vârful de memorie rezidentă raportat de runtime \\
\texttt{queue\_wait\_ms} & sistem & timpul dintre crearea rulării și preluarea job-ului de către runner \\
\texttt{result\_rows} & bază de date & numărul de puncte temporale persistate pentru rulare \\
\texttt{compartment\_count} & model & dimensiunea vectorului de stare salvat la fiecare pas \\
\texttt{api\_latency\_ms} & API & timpul de răspuns pentru crearea unei simulări și citirea rezultatului \\
\bottomrule
\end{tabularx}
\end{table*}

Primele trei metrici (\texttt{wall\_time\_ms}, \texttt{cpu\_time\_ms},
\texttt{peak\_memory\_mb}) sunt scrise de fiecare runner în tabela
\texttt{benchmarks} și sunt disponibile pentru toate rulările din campania
descrisă în secțiunea~\ref{subsec:rezultate_masurate}. Metricile derivate
(\texttt{queue\_wait\_ms}, latențe API, dimensiunea răspunsului JSON) ar
necesita o instrumentare suplimentară și nu fac parte din rezultatele
actuale; ele apar în Tabelul~\ref{tab:metrici_benchmark} pentru completitudine
a cadrului de evaluare.

\section{Matricea experimentală}
\label{sec:matrice_experimentala}

O evaluare completă trebuie să acopere toate combinațiile relevante dintre model,
limbaj și orizont temporal. Matricea minimă recomandată este prezentată în
Tabelul~\ref{tab:matrice_benchmark}.

\begin{table*}[!htbp]
\centering
\caption{Matrice minimă pentru campania de benchmark.}
\label{tab:matrice_benchmark}
\begin{tabularx}{0.95\textwidth}{l l X X}
\toprule
\textbf{Dimensiune} & \textbf{Valori} & \textbf{Scop} & \textbf{Observații} \\
\midrule
Model & SIS, SIR, SEIR, SIRS, SIDARTHE & cost în funcție de complexitate & două până la opt compartimente \\
Limbaj & R, Julia, Python & comparație între ecosisteme & prima rulare Julia trebuie separată de rulările calde \\
Orizont $t_{\max}$ & 365, 2\,000, 10\,000, 50\,000 zile & scalare cu numărul de pași de integrare și de puncte salvate & singurul lever real pentru rulările ODE; lungimea setului de date nu afectează durata solverului \\
Repetări & minim 10 & robustețe statistică & raportare prin mediană și IQR \\
Concurență & 1, 3, 5 job-uri & comportament operațional & măsoară coadă, DB și scheduling \\
\bottomrule
\end{tabularx}
\end{table*}

Această matrice permite două tipuri de întrebări. Prima este numerică: cât crește
durata când modelul trece de la două compartimente la opt? A doua este
operațională: cum se schimbă timpul total când mai multe simulări sunt lansate
simultan?

Două precizări sunt importante pentru interpretarea matricei. (i) Pentru
runnerele compartimentale (R, Julia, Python), durata unei rulări nu depinde de
mărimea setului de date --- runnerul citește doar prima înregistrare pentru
condițiile inițiale și integrează $t_{\max}$ pași. Levierul real pentru rulări
lungi este, prin urmare, $t_{\max}$, complexitatea modelului și toleranțele
solverului. (ii) Pentru runnerul LSTM, dimpotrivă, costul de antrenare scalează
cu lungimea seriei (numărul de eșantioane glisante), cu numărul de epoci și cu
dimensiunea stării ascunse; în această privință, seturile sintetice de zeci de
mii de zile sunt esențiale pentru evaluarea regimului saturat al rețelei.

Execuția campaniei este automatizată prin scriptul
\texttt{scripts/run\_benchmark.py}, care primește credențialele super-admin,
identificatorii unui studiu și ai unuia sau mai multor seturi de date, lista
modelelor, lista limbajelor, valorile $t_{\max}$ de baleiat și numărul de
repetări. Driverul lansează combinațiile prin API
(\texttt{POST /api/studies/\{id\}/simulations}), așteaptă finalizarea fiecărei
rulări prin polling și măsoară latența end-to-end pe partea de client (de la
crearea simulării până la statusul \texttt{completed}), care include coada
Redis, solverul și persistarea --- exact descompunerea propusă mai sus.
Rezultatele sunt scrise într-un CSV per campanie, iar metricile interne ale
solverului (\texttt{wall\_time\_ms}, \texttt{cpu\_time\_ms},
\texttt{peak\_memory\_mb}) rămân disponibile în tabela \texttt{benchmarks},
agregabile prin \texttt{GET /api/studies/\{id\}/benchmarks}.

\begin{table*}[!htbp]
\centering
\caption{Teste necesare pentru obținerea rezultatelor de evaluare software.}
\label{tab:teste_necesare_software}
\begin{tabularx}{0.96\textwidth}{lX X}
\toprule
\textbf{Zonă evaluată} & \textbf{Test necesar} & \textbf{Rezultat de raportat} \\
\midrule
Runner numeric & rulări repetate pentru fiecare model, limbaj și orizont temporal & mediană, IQR, minim și maxim pentru timp și memorie \\
Cold-start/Warm-start & separarea primei rulări de rulările ulterioare & costul de inițializare al runtime-ului, în special pentru Julia \\
Concurență & lansarea simultană a 1, 3 și 5 job-uri & timp total, timp de așteptare în coadă și degradarea latenței \\
Persistență rezultate & rulări cu 30, 90, 180 și 365 puncte temporale & timp de scriere, număr de rânduri și dimensiunea JSONB \\
API și frontend & citirea rezultatelor și exportul CSV pentru fiecare model & latență endpoint, dimensiune răspuns și timp de export \\
\bottomrule
\end{tabularx}
\end{table*}

\section{Protocol de rulare}
\label{subsec:protocoale}

Pentru rezultate reproductibile, campania de benchmark trebuie să respecte un
protocol fix:
\begin{enumerate}
    \item se pornește aceeași versiune de cod și aceeași configurație Docker;
    \item se rulează migrațiile bazei de date și se curăță datele experimentale
    vechi;
    \item se execută rulări de încălzire pentru fiecare limbaj și model;
    \item se rulează combinațiile din matricea experimentală;
    \item se exportă tabela \texttt{benchmarks} și metadatele din
    \texttt{simulation\_runs};
    \item se agregă rezultatele prin mediană, IQR, minim și maxim;
    \item se păstrează separat rezultatele cold-start și warm-start.
\end{enumerate}

Separarea cold-start/warm-start este necesară în special pentru Julia, unde prima
rulare poate include cost de compilare JIT. Pentru Python și R, costurile mari
apar de obicei la pornirea containerului și la importul bibliotecilor, nu la
fiecare job.

Campania de bază a fost executată pe clusterul \texttt{kube-4kdhb}
(3 worker-i $\times$ 4~vCPU / 8~GiB pe OpenStack), folosind imaginile
\texttt{emedicine-*:latest} de pe GHCR, construite din branch-ul principal. Rezultatele sunt prezentate în
secțiunea~\ref{subsec:rezultate_masurate}; reproducerea pe alte configurații
hardware (de exemplu o mașină locală sau un nod cu GPU) este planificată ca
script separat, în afara aplicației, și permite plasarea acelorași combinații
\emph{model$\,\times\,$limbaj$\,\times\,t_{\max}$} pe scări de hardware
diferite.

\section{Scalabilitate și complexitatea modelului}
\label{sec:scalabilitate_date}

În modelele compartimentale agregate, mărimea populației $N$ nu crește direct
dimensiunea sistemului diferențial. O simulare SIDARTHE pentru $N=10^3$ și una
pentru $N=10^7$ au același număr de ecuații; se schimbă doar scara valorilor.
Costul crește în schimb cu numărul de compartimente, cu orizontul temporal
$t_{\max}$, cu numărul de puncte salvate și cu numărul de fluxuri calculate în
funcția RHS.

SIDARTHE este mai costisitor decât SIR/SIS/SEIR/SIRS nu pentru că populația este
mai mare, ci pentru că fiecare evaluare RHS calculează mai multe tranziții
clinice și fiecare punct salvat conține mai multe valori. În plus, persistarea în
JSONB introduce un cost de serializare care nu exista într-o schemă rigidă
dedicată unui singur model, dar acest cost cumpără flexibilitate: aceeași tabelă
poate stoca modele cu structuri diferite de stare.

\section{Compararea implementărilor R, Julia și Python}
\label{sec:comparatie_timp}

Cele trei implementări nu folosesc exact același algoritm numeric. Runner-ul R
folosește RK4 cu pas fix $h=1$, Python folosește \texttt{solve\_ivp} cu RK45, iar
Julia folosește \texttt{Tsit5}. Comparația trebuie interpretată ca o comparație
între implementări practice echivalente funcțional, nu ca un experiment izolat pe
același solver. Pentru o comparație strictă ar fi necesară implementarea manuală a
aceleiași metode în toate cele trei limbaje.

Avantajul abordării curente este că reflectă utilizarea naturală a fiecărui
ecosistem: R prin \texttt{deSolve}, Python prin SciPy, Julia prin SciML.
Rezultatele benchmark-ului trebuie deci citite ca o comparație între ecosisteme de
calcul științific integrate într-un sistem software, nu doar între limbaje de
programare.

\section{Indicatori agregați}
\label{sec:indicatori_agregati}

Pentru fiecare combinație model--limbaj se raportează mediana timpului wall-clock:
\begin{equation}
\widetilde{T}_{m,l} = \operatorname{median}
\left(T_{m,l}^{(1)}, T_{m,l}^{(2)}, \ldots, T_{m,l}^{(r)}\right),
\end{equation}
unde $r$ este numărul de repetări. Variabilitatea este exprimată prin intervalul
intercuartilic:
\begin{equation}
\operatorname{IQR}_{m,l}=Q_{3,m,l}-Q_{1,m,l}.
\end{equation}

Pentru comparații între limbaje se poate raporta raportul de viteză față de
Python:
\begin{equation}
\operatorname{speedup}_{l} =
\frac{\widetilde{T}_{m,\text{Python}}}{\widetilde{T}_{m,l}}.
\end{equation}
Această valoare trebuie interpretată per model și per orizont temporal. Un limbaj
poate fi mai rapid pentru SIDARTHE pe 365 de zile și mai lent pentru SIR pe 30 de
zile dacă overheadul runtime-ului domină calculul numeric.

\subsection{Rezultate măsurate}
\label{subsec:rezultate_masurate}

A fost executată o campanie completă prin scriptul
\texttt{run\_benchmark.py}, conform protocolului din
secțiunea~\ref{subsec:protocoale}: cinci modele compartimentale, trei limbaje,
patru orizonturi temporale ($t_{\max} \in \{90, 365, 2\,000, 10\,000\}$),
3 repetări per combinație, pe setul de date România OWID. Toate cele 180 de
rulări s-au finalizat fără erori. Datele sunt extrase direct din tabela
\texttt{benchmarks} a aplicației deployate, valorile reflectând timpul măsurat
de către runner între începutul integrării ODE și finalul scrierii rezultatelor.

\begin{table*}[!htbp]
\centering
\caption{Wall-time median al solverului (ms), pe model si limbaj, pentru patru orizonturi $t_{\max}$. Rezultate din campania de benchmark prin aplicatie (cluster \texttt{kube-4kdhb}, 3 worker x 4 vCPU / 8 GiB), 3 repetari per combinatie, dataset Romania OWID, prima inregistrare ca conditie initiala.}
\label{tab:bench_wall_time}
\small
\begin{tabularx}{\textwidth}{l l rrrr}
\toprule
\textbf{Model} & \textbf{Limbaj} & \multicolumn{4}{c}{\textbf{Wall-time median (ms)}} \\
\cmidrule(l){3-6}
 & & $t_{\max}=90$ & $t_{\max}=365$ & $t_{\max}=2000$ & $t_{\max}=10000$ \\
\midrule
\texttt{sis} & \texttt{r} & 12.0 & 41.0 & 149.0 & 738.0 \\
\texttt{sis} & \texttt{julia} & 0.0 & 0.0 & 1.0 & 5.0 \\
\texttt{sis} & \texttt{python} & 4.0 & 5.0 & 4.0 & 6.0 \\
\midrule
\texttt{sir} & \texttt{r} & 12.0 & 35.0 & 146.0 & 807.0 \\
\texttt{sir} & \texttt{julia} & 0.0 & 1.0 & 1.0 & 5.0 \\
\texttt{sir} & \texttt{python} & 4.0 & 4.0 & 5.0 & 5.0 \\
\midrule
\texttt{seir} & \texttt{r} & 19.0 & 35.0 & 161.0 & 673.0 \\
\texttt{seir} & \texttt{julia} & 1.0 & 0.0 & 1.0 & 5.0 \\
\texttt{seir} & \texttt{python} & 4.0 & 3.0 & 5.0 & 9.0 \\
\midrule
\texttt{sirs} & \texttt{r} & 10.0 & 40.0 & 154.0 & 665.0 \\
\texttt{sirs} & \texttt{julia} & 0.0 & 1.0 & 1.0 & 5.0 \\
\texttt{sirs} & \texttt{python} & 5.0 & 4.0 & 5.0 & 5.0 \\
\midrule
\texttt{sidarthe} & \texttt{r} & 16.0 & 62.0 & 274.0 & 1327.0 \\
\texttt{sidarthe} & \texttt{julia} & 2.0 & 1.0 & 2.0 & 6.0 \\
\texttt{sidarthe} & \texttt{python} & 5.0 & 5.0 & 5.0 & 8.0 \\
\bottomrule
\end{tabularx}
\end{table*}

\begin{table*}[!htbp]
\centering
\caption{Raport de viteza al solverului fata de Python (mediana wall-time Python / mediana limbaj). Valori $>1$ inseamna mai rapid decat Python; valori $<1$ inseamna mai lent. Pragul `$>x$' apare cand mediana sub-milisecundara face raportul exact instabil.}
\label{tab:bench_speedup_vs_python}
\small
\begin{tabularx}{\textwidth}{l l rrrr}
\toprule
\textbf{Model} & \textbf{Limbaj} & \multicolumn{4}{c}{\textbf{Speedup vs Python}} \\
\cmidrule(l){3-6}
 & & $t_{\max}=90$ & $t_{\max}=365$ & $t_{\max}=2000$ & $t_{\max}=10000$ \\
\midrule
\texttt{sis} & \texttt{r} & 0.33 & 0.12 & 0.03 & 0.01 \\
\texttt{sis} & \texttt{julia} & $>8$ & $>10$ & 4.00 & 1.20 \\
\midrule
\texttt{sir} & \texttt{r} & 0.33 & 0.11 & 0.03 & 0.01 \\
\texttt{sir} & \texttt{julia} & $>8$ & 4.00 & 5.00 & 1.00 \\
\midrule
\texttt{seir} & \texttt{r} & 0.21 & 0.09 & 0.03 & 0.01 \\
\texttt{seir} & \texttt{julia} & 4.00 & $>6$ & 5.00 & 1.80 \\
\midrule
\texttt{sirs} & \texttt{r} & 0.50 & 0.10 & 0.03 & 0.01 \\
\texttt{sirs} & \texttt{julia} & $>10$ & 4.00 & 5.00 & 1.00 \\
\midrule
\texttt{sidarthe} & \texttt{r} & 0.31 & 0.08 & 0.02 & 0.01 \\
\texttt{sidarthe} & \texttt{julia} & 2.50 & 5.00 & 2.50 & 1.33 \\
\bottomrule
\end{tabularx}
\end{table*}

\begin{table}[!htbp]
\centering
\caption[Memorie peak (mediană)]{Memorie reziduală peak (MB), mediana pe toate combinațiile model x $t_{\max}$.}
\label{tab:bench_memory}
\footnotesize
\setlength{\tabcolsep}{4pt}
\begin{tabularx}{\columnwidth}{l X r}
\toprule
\textbf{Limbaj} & \textbf{Runtime} & \textbf{Memorie (MB)} \\
\midrule
\texttt{r} & deSolve + RPostgres & 129 \\
\texttt{julia} & DifferentialEquations.jl (JIT) & 940 \\
\texttt{python} & NumPy + SciPy & 95 \\
\bottomrule
\end{tabularx}
\end{table}


Vizual, scalarea wall-time-ului în funcție de $t_{\max}$ (pentru modelul cel
mai complex, SIDARTHE) este prezentată în
Figura~\ref{fig:scaling_wall_time}; cele trei limbaje produc trei regimuri
calitativ diferite, identificabile dintr-o privire pe o scară log-log.

\begin{figure}[t]
\centering
\begin{tikzpicture}
\begin{loglogaxis}[
    width=\columnwidth,
    height=6.5cm,
    xlabel={$t_{\max}$ (zile)},
    ylabel={Wall-time median al solverului (ms)},
    xmin=80, xmax=12000,
    ymin=0.5, ymax=2000,
    grid=both, grid style={line width=.1pt, draw=gray!20},
    minor tick num=9,
    legend style={at={(0.02, 0.98)}, anchor=north west, font=\footnotesize, fill=white, draw=gray!50},
    legend cell align={left},
    tick label style={font=\footnotesize},
    label style={font=\small},
    log basis x=10,
    log basis y=10,
]
\addplot[mark=square*, thick, color=blue!80!black, mark size=2.5pt]
    table[col sep=comma, x=t_max, y=julia] {data/scaling_wall_sidarthe.csv};
\addlegendentry{Julia (\texttt{Tsit5})}

\addplot[mark=triangle*, thick, color=green!50!black, mark size=2.5pt]
    table[col sep=comma, x=t_max, y=python] {data/scaling_wall_sidarthe.csv};
\addlegendentry{Python (\texttt{solve\_ivp} RK45)}

\addplot[mark=*, thick, color=red!80!black, mark size=2.5pt]
    table[col sep=comma, x=t_max, y=r] {data/scaling_wall_sidarthe.csv};
\addlegendentry{R (\texttt{deSolve} RK4)}

\end{loglogaxis}
\end{tikzpicture}
\caption{Scalarea wall-time-ului median al solverului SIDARTHE in functie de
orizontul $t_{\max}$, pentru cele trei limbaje. R prezinta o crestere
aproape liniara cu $t_{\max}$ ($\sim$$130\times$ intre 90 si 10\,000 de pasi),
dominata de inserarile rand-cu-rand in baza de date prin
\texttt{dbExecute()}; integrarea propriu-zisa cu \texttt{deSolve::ode} este
sub-secundara. Python ramane stabil ($4$--$8$~ms), iar Julia este mai rapid
in regimul amortizat dupa compilarea JIT, atingand $6$~ms la $10\,000$ de pasi.}
\label{fig:scaling_wall_time}
\end{figure}


Din date reies trei tipare dominante.

Julia este cel mai rapid solver, dar are un cost de pornire JIT
semnificativ. La $t_{\max}=10\,000$, mediana wall-time-ului Julia este de
$5$--$6$~ms pentru orice model, inclusiv SIDARTHE cu opt compartimente.
Python rămâne în intervalul $5$--$9$~ms, iar R urcă de la $665$ la $1\,327$~ms,
de două până la trei ordine de mărime peste celelalte. Pentru orizonturi mici
($t_{\max}=90$), prima rulare Julia plătește însă costul compilării
\emph{just-in-time}, vizibil în diferența dintre repetarea 1 și repetările 2--3
(Tabelul~\ref{tab:bench_julia_coldstart}): $\sim$$900$--$1\,500$~ms doar
pentru încărcarea și compilarea funcțiilor RHS și a solverului
\texttt{Tsit5}.

\begin{table}[!htbp]
\centering
\caption[Cost cold-start JIT Julia]{Costul de pornire JIT pentru runner-ul Julia: prima rulare include compilarea, repetările ulterioare rulează din cod cache-uit. Valori la $t_{\max}=90$.}
\label{tab:bench_julia_coldstart}
\resizebox{\columnwidth}{!}{%
\begin{tabular}{l r r r}
\toprule
\textbf{Model} & \textbf{Rep 1 (ms)} & \textbf{Mediana rep 2-3 (ms)} & \textbf{Cost JIT (ms)} \\
\midrule
\texttt{sis} & 1435 & 0 & 1435 \\
\texttt{sir} & 891 & 0 & 891 \\
\texttt{seir} & 887 & 0 & 886 \\
\texttt{sirs} & 948 & 0 & 948 \\
\texttt{sidarthe} & 2 & 1 & 1 \\
\bottomrule
\end{tabular}%
}
\end{table}


Python (SciPy \texttt{solve\_ivp}) este cel mai stabil ca performanță:
wall-time-ul rămâne aproape constant în jurul a $4$--$9$~ms indiferent de
model sau orizont, fără cold-start vizibil. Acest comportament îl face
limbajul implicit recomandabil pentru utilizare interactivă, unde
predictibilitatea contează la fel de mult ca viteza brută.

Costul mare al runner-ului R la orizonturi mari nu vine de la \texttt{deSolve},
ci de la inserările individuale în baza de date. Cele $10\,001$ inserții
rând-cu-rând prin \texttt{dbExecute()} domină cele $\sim$$1$ secundă măsurată
la $\texttt{sidarthe}/\texttt{r}/t_{\max}=10\,000$; integrarea numerică
propriu-zisă cu \texttt{deSolve::ode} este sub-secundară. Avem astfel un
punct concret de îmbunătățire ulterioară (de exemplu prin \texttt{COPY FROM}
sau prin tranzacții batch), nu o limitare intrinsecă a ecosistemului R.

În ce privește memoria reziduală peak (Tabelul~\ref{tab:bench_memory}), Julia
plătește un overhead constant de $\sim$$940$~MB pentru runtime-ul SciML, în
timp ce Python rămâne sub $100$~MB iar R sub $135$~MB. Acest cost este
independent de complexitatea modelului și de orizont, deci nu se amortizează
la rulări scurte; pentru job-uri foarte multe și mici, Python și R sunt
opțiunile mai economicoase.

\subsection{Costul componentei LSTM}
\label{subsec:bench_lstm}

Runner-ul LSTM are un profil de cost diferit calitativ de runnerele
compartimentale. Costul ODE depinde de orizontul de integrare $t_{\max}$
(număr de pași și de puncte salvate), în timp ce costul LSTM este dominat de
\emph{antrenare} --- numărul de epoci $\times$ numărul de eșantioane glisante
--- iar $t_{\max}$ controlează doar orizontul prognozei autoregresive
(neglijabil în comparație).

Tabelul~\ref{tab:bench_lstm} arată rezultatele pe setul România OWID
(1\,440 zile, $80\%$ antrenare $/$ $20\%$ test cronologic) cu parametrii
impliciti din API. Mediana wall-time-ului este $\sim$$290$~s, aproape identică
între $t_{\max}=30$ și $t_{\max}=90$, ceea ce confirmă empiric că forecast-ul
nu contribuie semnificativ. Utilizarea CPU este peste $98\%$ din wall-time
(LSTM hand-rolled în NumPy fără paralelism), iar memoria peak este de doar
$\sim$$56$~MB, mai mică decât a runnerelor ODE în Python --- nu se construiesc
matrice de stare comparabile ca dimensiune.

\begin{table}[!htbp]
\centering
\caption{Costul rularii pentru runner-ul LSTM dedicat (Python, hand-rolled NumPy), parametri impliciti (\texttt{lookback=14}, \texttt{hidden\_size=12}, \texttt{epochs=150}, \texttt{learning\_rate=0.01}), dataset Romania OWID (1\,440 zile, $80\%$ antrenare / $20\%$ test cronologic). Forecast-ul autoregresiv pe orizontul $t_{\max}$ are cost neglijabil fata de antrenare.}
\label{tab:bench_lstm}
\begin{tabularx}{\columnwidth}{r r r r r}
\toprule
$t_{\max}$ & \textbf{n} & \textbf{Wall median (s)} & \textbf{CPU median (s)} & \textbf{Memorie (MB)} \\
\midrule
30 & 2 & 296.8 & 295.3 & 55.6 \\
90 & 2 & 286.2 & 285.0 & 56.3 \\
\bottomrule
\end{tabularx}
\end{table}


Pentru contextul tezei, această valoare absolută este mai puțin importantă
decât observația de scalare: pe seriile sintetice de $20\,000$ de zile,
antrenarea ar crește la $\sim$$70$ minute cu aceiași parametri, ceea ce
plasează LSTM ca operație asincronă lungă --- exact rolul pentru care
\texttt{runner-lstm} este izolat într-o coadă Redis dedicată.

\subsection{Variabilitatea cross-hardware}
\label{subsec:hardware_variability}

Toate rezultatele de mai sus provin de pe un singur profil hardware
(clusterul \texttt{kube-4kdhb}). Pentru a separa contribuția hardware-ului
de cea a stack-ului software, aceeași logică de solver a fost rulată ca
script standalone (\texttt{scripts/bench\_local\_solver.py}) pe o mașină
locală complet diferită --- Apple~M1~Pro arm64, nativ, fără container și
fără persistare în baza de date --- folosind exact aceleași tolerante numerice
($\texttt{RK45}$, $\mathrm{rtol}=10^{-8}$, $\mathrm{atol}=10^{-10}$), aceiași
parametri și aceleași condiții inițiale ca rulările prin aplicație. Așa cum
este indicat explicit în secțiunea~\ref{subsec:protocoale}, această
masurare se face în afara aplicației pentru a evita amestecarea variațiilor
de hardware cu deciziile de ortografie ale containerului și ale bazei de
date distribuite.

Tabelul~\ref{tab:hardware_comparison} compară median wall-time pentru SIR și
SIDARTHE la trei orizonturi temporale, pe workload-ul sintetic non-trivial
(setul \texttt{synthetic\_multiwave}, $I_0 = 100$ pentru a forța solverul să
parcurgă efectiv pași de integrare). Datele de cluster sunt re-rulate în
mod identic prin API (între $1$ și $3$ repetări per combinație, în funcție
de campania anterioară), iar datele M1 provin din scriptul local
(3 repetări per combinație, cu o rulare de încălzire ignorată).

\begin{table*}[!htbp]
\centering
\caption{Comparatia timpilor de executie pentru solverul Python (\texttt{scipy.solve\_ivp} RK45) pe doua profile hardware distincte: cluster \texttt{kube-4kdhb} (worker amd64, OpenStack VM, 4 vCPU / 8 GiB, runner in container, wall-time include si scrierea rezultatelor in YugabyteDB) vs M1 Pro arm64 nativ (script standalone, fara container, fara DB). Workload non-trivial: set sintetic \texttt{multiwave} cu $I_0 = 100$, $N = 5 \cdot 10^7$. Coloana \emph{speedup M1} este raportul medianelor; valori $>1$ inseamna M1 mai rapid.}
\label{tab:hardware_comparison}
\begin{tabularx}{\textwidth}{l r r r r r r}
\toprule
\textbf{Model} & \textbf{$t_{\max}$} & \textbf{$n_{\mathrm{cluster}}$} & \textbf{Cluster wall (ms)} & \textbf{$n_{\mathrm{M1}}$} & \textbf{M1 wall (ms)} & \textbf{Speedup M1} \\
\midrule
\texttt{sidarthe} & 365 & 1 & 86.0 & 3 & 16.44 & 5.2$\times$ \\
 & 2000 & 3 & 127.0 & 3 & 31.46 & 4.0$\times$ \\
 & 10000 & 3 & 345.0 & 3 & 90.48 & 3.8$\times$ \\
\texttt{sir} & 365 & 1 & 73.0 & 3 & 12.25 & 6.0$\times$ \\
 & 2000 & 3 & 96.0 & 3 & 16.07 & 6.0$\times$ \\
 & 10000 & 3 & 100.0 & 3 & 22.98 & 4.4$\times$ \\
\bottomrule
\end{tabularx}
\end{table*}


Pentru același cod Python și aceeași încărcare numerică, M1 Pro nativ
este de $\sim$$4\times$--$6\times$ mai rapid decât worker-ul amd64 al
clusterului. Acest raport include diferența de microarhitectură (Apple
Silicon vs.\ CPU x86 virtualizat în KVM), diferența de mediu (nativ vs.\
container) și costul redus al stratului de persistare (M1 nu scrie
rezultatele într-o bază de date distribuită cu factor de replicare 3).
Raportul este relativ stabil între modele și orizonturi, deci marea parte
a diferenței vine din execuția solver-ului propriu-zis, nu din
overhead-uri fixe.

Pentru workload-ul trivial (setul România OWID cu $I_0 = 0$, ca în
campania prin aplicație), M1 nativ rulează în $0{,}4$--$0{,}8$~ms per
solver, iar cluster-ul în $3$--$9$~ms (Tabelul~\ref{tab:bench_wall_time});
diferența crește la $10\times$ sau mai mult tocmai pentru că, atunci când
solver-ul nu are dinamică reală, \emph{toată} cheltuiala wall-time pe
cluster vine din scrierea celor $t_{\max}+1$ rânduri în
\texttt{simulation\_results}, cost care nu există în scriptul local.

Concluzia metodologică este că un raport \emph{absolut} între cele două
hardware-uri depinde semnificativ de natura workload-ului (solver-light vs.
solver-heavy) și de stack-ul software inclus în măsurare. Recomandarea de
limbaj din secțiunea~\ref{sec:analiza_sintetica} rămâne valabilă pe
clusterul testat; pentru o decizie cantitativă pe un alt profil
hardware --- în special cu accelerare GPU pentru LSTM sau cu solvere
implementate în Fortran/C++ legate static --- aceeași campanie trebuie
re-rulată local, iar scriptul \texttt{bench\_local\_solver.py} este punctul
de pornire pentru această replicare.

\section{Acuratețe numerică și verificări de regresie}
\label{sec:acuratete}

Corectitudinea rezultatelor este verificată prin conservarea populației. Pentru
SIDARTHE, suma $S+I+D+A+R+T+H+E$ trebuie să rămână egală cu $N$ la fiecare pas
salvat; pentru SEIR, $S+E+I+R=N$; iar pentru modelele SIR/SIS/SIRS se verifică
suma compartimentelor definite de fiecare model.

Un al doilea criteriu este apropierea traiectoriilor obținute în cele trei
limbaje pentru aceeași configurație. Deoarece solverele diferă, egalitatea exactă
nu este realistă. Se verifică în schimb abaterea maximă normalizată:
\begin{equation}
\Delta_{l_1,l_2} =
\max_{t,c}
\frac{\left|x_{l_1,c}(t)-x_{l_2,c}(t)\right|}{N}.
\end{equation}
Dacă această abatere rămâne sub un prag ales, implementările sunt considerate
compatibile funcțional pentru scopul aplicației.

Verificarea a fost rulată concret pe setul de date sintetic
\texttt{multiwave} (populație $N = 50 \cdot 10^6$), pentru care condițiile
inițiale produc dinamică ne-trivială (spre deosebire de România OWID, unde
prima zi din serie are $I=0$). Rezultatele sunt sintetizate în
Tabelul~\ref{tab:cross_language_regression}.

\begin{table}[!htbp]
\centering
\caption[Regresie numerică cross-limbaj]{Verificare numerică de regresie cross-limbaj: abaterea maximă absolută între traiectoriile compartimentale produse în R, Julia și Python pentru aceiași parametri și aceleași condiții initiale, pe setul sintetic \texttt{multiwave} (populația $N=50 \cdot 10^6$, $t_{\max}=365$ pași zilnici). Modelele folosesc solvere diferite (R cu RK4 pas fix $h=1$, Python cu RK45 adaptiv, Julia cu Tsit5), deci abaterile reflectă toleranțele numerice cumulate, nu erori de implementare.}
\label{tab:cross_language_regression}
\resizebox{\columnwidth}{!}{%
\begin{tabular}{l l r r}
\toprule
\textbf{Model} & \textbf{Pereche} & \textbf{Max $|\Delta|$} & \textbf{Normalizat ($/N$)} \\
\midrule
\texttt{sir} & \texttt{julia} vs \texttt{python} & 0.5 & $1{,}01 \cdot 10^{-8}$ \\
 & \texttt{julia} vs \texttt{r} & 2185.8 & $4{,}37 \cdot 10^{-5}$ \\
 & \texttt{python} vs \texttt{r} & 2186.2 & $4{,}37 \cdot 10^{-5}$ \\
\midrule
\texttt{sidarthe} & \texttt{julia} vs \texttt{python} & 0.2 & $3{,}72 \cdot 10^{-9}$ \\
 & \texttt{julia} vs \texttt{r} & 6570.9 & $1{,}31 \cdot 10^{-4}$ \\
 & \texttt{python} vs \texttt{r} & 6571.1 & $1{,}31 \cdot 10^{-4}$ \\
\bottomrule
\end{tabular}%
}
\end{table}


Implementările Julia și Python produc traiectorii practic identice
($\sim$$10^{-8} \cdot N$, la limita preciziei flotante), pentru că
ambele folosesc solvere Runge--Kutta adaptive cu toleranțe similare
($\texttt{Tsit5}$ și respectiv $\texttt{RK45}$). Implementarea R, care
folosește RK4 cu pas fix $h = 1$~zi, deviază de la celelalte două cu
$\sim$$10^{-4} \cdot N$ ($\sim$$2\,200$ cazuri pentru SIR, $\sim$$6\,600$
pentru SIDARTHE). Pentru o populație de $50$ de milioane, abaterea
reprezintă $0{,}013\%$ din populație, sub-vizuală pe o reprezentare
grafică tipică, compatibilă funcțional cu scopul aplicației, dar
deliberat mai puțin precisă decât metodele adaptive.

\section{Evaluarea API-ului și a bazei de date}
\label{sec:evaluare_api_db}

Performanța percepută de utilizator depinde și de API și baza de date. Pentru
crearea unei simulări, latența API include validarea inputului, citirea primului
rând epidemiologic, inserarea în \texttt{simulation\_runs} și publicarea în
Redis. Pentru vizualizarea rezultatului, latența include citirea rândurilor din
\texttt{simulation\_results}, deserializarea JSONB și construirea răspunsului.

Din acest motiv, benchmark-ul numeric a fost completat cu o campanie de teste de
API, realizată printr-un script reproductibil (\texttt{scripts/bench\_api.py}).
Pentru fiecare model s-au măsurat cele patru operații expuse de API și percepute
direct de utilizator: \emph{create} (crearea unei simulări), \emph{status}
(interogarea stării), \emph{results} (citirea traiectoriei complete în format
JSON) și \emph{export} (exportul CSV al aceleiași traiectorii). Latențele
\emph{create} au fost colectate din $12$ repetări, iar cele de citire din $30$ de
repetări per combinație; măsurătorile au fost făcute prin \texttt{port-forward}
către serviciul \texttt{api} din clusterul \texttt{kube-4kdhb}, deci includ tot
lanțul HTTP $\rightarrow$ API $\rightarrow$ YugabyteDB $\rightarrow$ HTTP, plus
variabilitatea de rețea introdusă de tunel.

\begin{table*}[!htbp]
\centering
\caption{Latenta API pe model, la orizontul de $\,365\,$ puncte (1 an, pas zilnic). Pentru fiecare endpoint se raporteaza mediana ($p_{50}$); pentru \emph{results}, endpointul cel mai greu, se raporteaza si percentila 95 ($p_{95}$) si dimensiunea raspunsului JSON. \emph{create} raspunde asincron (inainte de rularea solver-ului); \emph{results} si \emph{export} includ deserializarea din YugabyteDB. Masurat prin port-forward la clusterul \texttt{kube-4kdhb}.}
\label{tab:api_latency}
\small
\begin{tabularx}{\textwidth}{l r r r r r r}
\toprule
\textbf{Model} & \textbf{Create} & \textbf{Status} & \textbf{Results $p_{50}$} & \textbf{Results $p_{95}$} & \textbf{JSON} & \textbf{Export} \\
 & (ms) & (ms) & (ms) & (ms) & (KB) & (ms) \\
\midrule
\texttt{sir} & 46.4 & 23.9 & 48.5 & 139.8 & 38.1 & 48.3 \\
\texttt{sis} & 70.6 & 31.5 & 45.5 & 123.9 & 36.0 & 39.6 \\
\texttt{seir} & 49.9 & 27.0 & 53.8 & 155.6 & 40.3 & 36.6 \\
\texttt{sirs} & 50.0 & 28.3 & 46.9 & 124.0 & 38.1 & 37.3 \\
\texttt{sidarthe} & 44.2 & 29.7 & 47.1 & 126.3 & 48.9 & 39.3 \\
\bottomrule
\end{tabularx}
\end{table*}


Tabelul~\ref{tab:api_latency} prezintă rezultatele la orizontul de un an
($365$ puncte). Latențele mediane sunt omogene între modele și se situează
în intervalul $24$--$71$~ms pentru toate endpointurile; modelul nu
influențează semnificativ latența, deoarece costul este dominat de drumul
dus-întors prin baza de date, nu de structura compartimentală. Endpointul
\emph{create} răspunde asincron: doar validează inputul, citește primul
rând epidemiologic, inserează în \texttt{simulation\_runs} și publică în
coada Redis, înainte ca solverul să ruleze, deci latența sa nu reflectă
costul de integrare numerică, ci doar pe cel de inițiere a unui job.
Dimensiunea răspunsului JSON depinde de model: SIDARTHE, cu opt
compartimente, produce $\sim$$49$~KB față de $\sim$$36$~KB pentru SIS cu
două compartimente, la același orizont. Percentila 95 este de
$2{,}5$--$3$ ori mai mare decât mediana ($124$--$156$~ms), reflectând cozile
de latență tipice unei măsurători prin tunel; pe un acces direct în cluster
acestea ar fi mai mici, dar raportarea lor este corectă pentru profilul
de testare folosit.

\begin{table}[!htbp]
\centering
\caption[Scalare latență read cu orizontul]{Scalarea latenței de citire cu orizontul, pentru modelul \texttt{sidarthe} ($8$ compartimente). Pe măsură ce numărul de puncte crește, latența \emph{results} și dimensiunea JSON cresc aproximativ liniar, în timp ce \emph{export} CSV rămâne mai ieftin (text plat, fără structură JSON).}
\label{tab:api_latency_scaling}
\small
\resizebox{\columnwidth}{!}{%
\begin{tabular}{r r r r r}
\toprule
\textbf{Puncte} & \textbf{Results $p_{50}$} & \textbf{Results $p_{95}$} & \textbf{JSON} & \textbf{Export} \\
 & (ms) & (ms) & (KB) & (ms) \\
\midrule
30 & 30.5 & 135.6 & 4.1 & 31.1 \\
90 & 34.3 & 129.7 & 12.1 & 35.5 \\
180 & 40.0 & 113.4 & 24.1 & 34.7 \\
365 & 47.1 & 126.3 & 48.9 & 39.3 \\
\bottomrule
\end{tabular}%
}
\end{table}


Tabelul~\ref{tab:api_latency_scaling} izolează efectul orizontului pentru modelul
cel mai greu (SIDARTHE), variind numărul de puncte de la $30$ la $365$.
Dimensiunea răspunsului JSON crește aproape liniar cu numărul de puncte (de la
$\sim$$4$~KB la $\sim$$49$~KB), în timp ce latența mediană \emph{results} crește
mult mai lent (de la $\sim$$30$ la $\sim$$47$~ms): la aceste dimensiuni,
costul fix al interogării domină costul de serializare. Exportul CSV produce
constant răspunsuri de $5$--$10$ ori mai mici decât JSON-ul echivalent (text plat,
fără chei repetate per rând), confirmând că formatul de transport contează pentru
volumul transferat, chiar dacă latența rămâne comparabilă.

Concluzia practică este că, pentru orizonturile uzuale ($\le 365$ puncte) și
modelele studiate, limita sistemului nu este nici solverul (sub $10$~ms pe
cluster, conform Tabelului~\ref{tab:bench_wall_time}), nici transportul JSON, ci
costul fix al drumului prin baza de date, de ordinul a câtorva zeci de
milisecunde per cerere. Aceasta este o latență confortabilă pentru o interfață
interactivă și indică faptul că optimizările viitoare ar trebui să vizeze
reducerea numărului de cereri (de exemplu prin agregarea rezultatelor mai multor
rulări într-un singur răspuns), nu accelerarea unei singure cereri.

\section{Interpretarea rezultatelor}
\label{sec:analiza_sintetica}

Evaluarea corectă trebuie să combine trei axe: eroarea numerică, timpul de
execuție și complexitatea modelului. Un model nu este automat superior doar
pentru că are eroare mai mică. Dacă obține această eroare cu un număr mare de
parametri și cu timp de calibrare ridicat, poate fi mai puțin potrivit pentru
analize rapide. În sens invers, un model simplu și rapid poate fi foarte util
pentru explorare, dar insuficient pentru întrebări clinice.

Recomandarea practică derivată din această campanie urmează profilul cazului
de utilizare. Pentru rulări numeroase și scurte (de exemplu explorarea
sensibilității parametrilor), Python este alegerea implicită cea mai bună:
timpii sunt sub $10$~ms și predictibili pentru toate modelele și orizonturile
testate, fără cost de pornire. Pentru rulări puține și foarte lungi (de
exemplu integrare pe orizonturi $t_{\max} \gg 10\,000$ sau calibrare cu mii
de iterații în care Julia se compilează o singură dată), Julia oferă un
solver cu un ordin de mărime mai rapid, cu condiția ca cele $\sim$$1$~secundă
de pornire să fie amortizate. R rămâne potrivit pentru explorare interactivă
în pipeline-uri statistice familiare cercetătorilor din epidemiologie, însă
runner-ul curent are o limitare operațională (inserțiile rând-cu-rând în
baza de date) care domină wall-time-ul la orizonturi mari și poate fi
corectată independent de schimbarea limbajului.

Pentru predicție cu LSTM, costul este dominat de antrenare și nu se compară
direct cu solverele ODE. Este important ca interfața utilizatorului să trateze
LSTM ca operație asincronă explicită (lansare, polling, finalizare după
minute), așa cum este deja proiectată în aplicație.

Toate aceste afirmații sunt valide pentru profilul hardware al clusterului
\texttt{kube-4kdhb}. Reproducerea pe alte configurații (mașini locale, noduri
mai puternice, eventual GPU pentru LSTM) este lăsată ca extensie a campaniei,
prin scripturi rulate în afara aplicației.
